\documentclass{article}
	\usepackage[UTF8]{ctex}
	\usepackage{amsmath}
	\usepackage{graphicx}
	\usepackage{underscore}
	\usepackage{geometry}
	\usepackage{anysize}
	\date{}
	\title{时间差 \\ [2ex] \begin{large} 
	时间限制：C/C++ 1秒，其他语言2秒 \\
	空间限制：C/C++ 32.000MB，其他语言64.000MB 	\\
	64bit IO Format: \%lld
 			\end{large} }
	\geometry{top=10mm,bottom=15mm,left=25mm,right=25mm}
	
	\begin{document}
		\maketitle \vspace{-50pt}
		
		\section*{题目描述}
		\indent 给出同一天的两个时间，请求出两个时间的时间差。

		\section*{输入描述}
			\noindent 第1行输入''hh:mm:ss''格式的时间1。\\
			\noindent 第2行输入''hh:mm:ss''格式的时间2。\\
			\noindent hh代表时，mm代表分，ss代表秒。hh是[0,23]范围内的一个整数，mm是[0,59]范围内的一个整数，ss是[0,59]范围内的一个整数。
		\section*{输出描述}
			\noindent 输出''hh:mm:ss''格式的时间1与时间2的时间差。
			
		\section*{示例1}
			\subsection*{输入}
\noindent
01:02:03\\
04:05:06

			\subsection*{输出}
\noindent
03:03:03

		\section*{示例2}
			\subsection*{输入}
\noindent
01:02:03\\
00:05:06

			\subsection*{输出}
\noindent
00:56:57
			
		\section*{备注}

	\end
	{document}
	